% =====================================================================
%  Coupled UAV transportation and communication scheduling in mountain
%  flood disaster relief -- IEEE journal manuscript.
%  Compile with:  pdflatex main_en.tex
% =====================================================================
\documentclass[journal]{IEEEtran}
\IEEEoverridecommandlockouts

\usepackage{cite}
\usepackage{amsmath,amssymb,amsfonts}
\usepackage{algorithmic}
\usepackage{algorithm}
\usepackage{graphicx}
\usepackage{textcomp}
\usepackage{xcolor}
\usepackage{booktabs}
\usepackage{multirow}
\usepackage{flushend}
\allowdisplaybreaks
% compact float spacing (IEEE defaults are generous; keeps the paper within 10 pages)
\setlength{\textfloatsep}{6pt plus 1.0pt minus 2.0pt}
\setlength{\dbltextfloatsep}{8pt plus 1.0pt minus 2.0pt}
\setlength{\intextsep}{6pt plus 1.0pt minus 2.0pt}
\setlength{\abovecaptionskip}{3pt}
\setlength{\belowcaptionskip}{1pt}
\setlength{\abovedisplayskip}{5pt plus 1pt minus 1pt}
\setlength{\belowdisplayskip}{5pt plus 1pt minus 1pt}
\setlength{\abovedisplayshortskip}{3pt plus 1pt}
\setlength{\belowdisplayshortskip}{3pt plus 1pt}

\begin{document}

\title{A Unified Multi-Model Framework for Coupled UAV Transportation and
Communication Scheduling in Mountain Flood Disaster Relief}

\author{
\IEEEauthorblockN{Yang Xin\textsuperscript{*}, Zhao Zidong, Song Yuwei}
\IEEEauthorblockA{\textit{School of Mechanical and Power Engineering, Zhengzhou University}\\
Zhengzhou, Henan 450001, China\\
Email: xinyang.zzu@outlook.com}
\thanks{\textsuperscript{*}Corresponding author: Yang Xin
(e-mail: xinyang.zzu@outlook.com).}
}

\maketitle
\markboth{Manuscript for Review, October 2026}%
{Yang \MakeLowercase{\textit{et al.}}: A Unified Multi-Model Framework for Coupled UAV Transportation and Communication Scheduling}

\begin{abstract}
Mountain floods damage roads and communication facilities simultaneously, so
ground vehicles can hardly enter the affected areas and UAVs become a key means
of material delivery and communication support. Severe terrain relief, strict
task deadlines and the tight coupling between transportation and communication
call for both to be planned under one physical caliber. This paper builds a
``transport scheduling--communication support--resource allocation'' chain
solved in four sequential stages; its central stage minimizes the relay
resources that keep the delivery UAVs connected once their delivery plan is
fixed, while batching, routing and partitioning generate and assess that plan.
Problem 1 (DEM-ELBPP) is an energy-aware bin-packing model solved exactly by
set-partitioning dynamic programming, giving an 18-trip plan at 59.24~kWh that
saves 0.69\% of energy over a greedy FFD/BFD baseline on the same instance.
Problem 2 (HF-VRPBT) adds heterogeneous fleet routing with time windows and
recommends a 20-trip plan at 65.25~kWh satisfying all 31 hard deadlines.
Problem 3 (TD-CJO) introduces a \emph{multi-hover-point direct-transfer} relay
mode: two relay UAVs cover all 11683 sampled relay-required points (11615 unique
sortie--time pairs) on the 1-s grid (relay energy 5.43~kWh, joint 70.68~kWh),
whereas a single-hover baseline covers only 62.7\% with two UAVs and 82.1\% with
three, so the mode is what makes grid-level full coverage attainable.
Problem 4 (SC-MAE-RG) compares $K{=}2$ and $K{=}3$ partitions (31/38
units, gaps 5/10) and recommends $K{=}2$. An independent recomputation reproduces
the point count within 0.3\% and the relay energy within 1.6\%, and the weakest
relay link keeps 0.34~dB of margin over the $-90$~dBm threshold.
\end{abstract}

\begin{IEEEkeywords}
mountain flood, UAV transportation scheduling, communication relay deployment,
multi-model fusion, resource-gap accounting
\end{IEEEkeywords}

% =====================================================================
\section{Introduction}
% ---------------------------------------------------------------------
Mountain flood disasters strike suddenly and are spatially dispersed, often
destroying roads and communication facilities at the same time. Continuous heavy
rainfall may trigger flash floods, landslides and road collapses, isolating
residents from stable ground transportation for a short period; in the early
post-disaster stage, drinking water, emergency food, medical supplies and
hygiene items must be delivered as soon as possible, while conventional vehicles
can hardly enter the strongly undulating and severely damaged areas. UAVs are
less sensitive to ground-road conditions and can serve as an important
complementary means of emergency transportation.

In July 2026, torrential rainfall brought by typhoon ``Maysak'' caused severe
flooding in Zhenlong Township, Hengzhou, Guangxi. Field reports indicated that
flood water, rockfall and road collapse cut off access from several directions,
leaving the core area temporarily without road traffic and with unstable
communication; the township suffered simultaneous road, power and network
outages and about 8,000 residents were trapped. During the response, food,
drinking water and medicine had to be carried in on foot or dropped by air, and
large UAVs also transported power-repair materials, while temporary
communication and emergency power equipment were deployed to gradually restore
contact. This episode shows that emergency delivery and communication support
must be organized jointly under blocked roads, tight energy budgets and
incomplete information.

From a technical viewpoint, completing delivery depends not only on whether a
UAV can reach the target area, but also on payload, horizontal flight distance,
climb/descent height, along-route terrain clearance, energy reservation, task
deadlines and equipment turnover. Meanwhile, a transportation UAV must maintain
command, telemetry and hand-over links during climb, cruise, descent and
hand-over; when fixed communication facilities are damaged or terrain blocks the
ground gateway, airborne relay UAVs must provide temporary coverage at suitable
positions and time windows. Therefore, material delivery under mountain floods
involves transportation organization plus simultaneous communication support and
equipment-resource coordination --- a typical transportation--communication--
resource coupled optimization problem.

The four stages studied here are ordered by dependency rather than by
convention. Batching (Problem 1) fixes the indivisible unit loads and their
per-area energy cost before any timing decision is made, because a route
cannot be timed before its payload is known. Routing (Problem 2) then adds
deadlines, airframe concurrency and battery turnover on top of these fixed
batches. Relay scheduling (Problem 3) is meaningful only after the routing
timeline exists, since the relay demand is exactly the set of trajectory points
whose link to the gateway fails. Finally, resource partitioning (Problem 4)
accounts peak concurrency on the joint plan, so it must come last. Each stage
therefore consumes only the outputs of its predecessors, which is what allows
one physical caliber to run through the whole chain without re-fitting.

\subsection{Related Work}
% ---------------------------------------------------------------------
Relay deployment under terrain blockage has been shown to be critical for link
quality: \cite{cuu2026} showed that even a single well-placed relay can
significantly improve link reliability in UAV-assisted robust communication link
design. In emergency logistics, \cite{wu2024} built a transportation-path
optimization model revealing the necessity of a timeliness--energy trade-off.
Cellular-connected cargo delivery and joint UAV--ground routing have since been
studied extensively, including trajectory optimization for cellular-connected
delivery fleets \cite{cao2024}, joint delivery routing of UAVs and ground
vehicles from an urban-logistics perspective \cite{wuXu2026} and under regional
restrictions \cite{wu2026}, and multi-objective design of UAV transportation
networks \cite{xia2019}. In cooperative mountain logistics, \cite{sun2026}
planned drone--vehicle missions for orchard fruit transport and \cite{dong2026}
developed three-dimensional path planning and scheduling for urban emergency UAVs
under typhoon disasters, showing the value of routing models in mountainous and
post-disaster settings; related disaster-response work has also addressed the
design of dedicated flood search-and-rescue UAVs \cite{wan2024}. On the
communication side, \cite{sabuncu2024} optimized drone-to-everything links for
medical services and \cite{huang2024} improved UAV-assisted downlink
transmission, while \cite{zhang2021} studied self-triggered cooperative control
of tethered multi-UAV transportation under low communication, and \cite{li2026}
surveyed UAV--UGV collaborative sensing for low-altitude transportation.
Heterogeneous multi-UAV swarms \cite{ma2020} and two-stage UAV facility location
and scheduling with objective-robustness optimization \cite{guo2026} are also
closely related. These studies, however, treat transportation planning and
communication provisioning \emph{separately}: routing works rarely enforce
continuous link coverage, while communication works rarely account for payload,
deadline and battery-turnover limits. The two closest studies still leave the
coupling open: the objective-robust siting model of \cite{guo2026} decides
where facilities stand but not how airborne links are maintained during
delivery, and the UAV--UGV sensing survey of \cite{li2026} catalogs
collaboration patterns without a joint transport--communication optimizer. Two
further works couple the two sides more directly: \cite{yanmaz2023} quantifies
the penalty of decoupled versus joint multi-UAV path planning for search and
rescue, and \cite{soya2025} designs delivery-UAV routes that reserve multi-hop
relays. Neither, however, carries payload, deadline or turnover constraints on a
common caliber. The gap addressed here is therefore the
joint, physically consistent modeling of transport batching, relay deployment and
resource turnover inside one optimization chain. Table~\ref{tab:rw} positions the
above studies against the three capabilities that the coupled problem requires
simultaneously: continuous link coverage (L), payload/energy feasibility (P) and
battery/component turnover (T).

\begin{table}[t]
\centering
\footnotesize
\setlength{\tabcolsep}{3pt}
\caption{Positioning of the related work. L = continuous link coverage,
P = payload/deadline feasibility, T = battery/component turnover.}
\label{tab:rw}
\begin{tabular}{llccc}
\toprule
Stream & Refs. & L & P & T \\
\midrule
Transport \& delivery routing &
\cite{wu2024,cao2024,wuXu2026,wu2026,xia2019,sun2026,dong2026,wan2024}
& -- & \checkmark & -- \\
UAV communication \& relay & \cite{cuu2026,sabuncu2024,huang2024,yanmaz2023,soya2025}
& \checkmark & -- & -- \\
Platform, swarm \& control & \cite{zhang2021,ma2020}
& -- & -- & \checkmark \\
Siting, sensing \& surveys & \cite{li2026,guo2026}
& -- & \checkmark & -- \\
This work & & \checkmark & \checkmark & \checkmark \\
\bottomrule
\end{tabular}
\end{table}

\subsection{Contributions}
The main contributions are:
\begin{itemize}
\item \textbf{A unified, reproducible model chain.} We build a
``transport scheduling--communication support--resource allocation'' framework in
which a single physical caliber (DEM-based cruise altitude, nonlinear equivalent
range, climb additional energy and two-stage charging) runs through all four
stages, so that results of successive stages are directly reusable.
\item \textbf{Multi-hover-point relay scheduling (MHPRS).} We formalize a relay
mode in which one relay UAV serves several hover points within a single
take-off/return cycle. Two relay UAVs then reach full mission coverage
(11683/11683 sampled relay points), whereas a single-hover-point round-trip
baseline covers only 62.7\% with two UAVs and 82.1\% with three --- even three
single-hover relays leave 17.9\% of the sampled points unserved, so the transfer
between hover points is the decisive factor. The formal definition and a
mode-dominance proposition are given in Sec.~\ref{sec:mhprs}.
\item \textbf{DEM-driven energy-aware bin packing.} We discretize continuous
energy constraints into bin-packing feasibility, enabling exact batching (18
trips) and a verified optimality benchmark. On the real instance the exact
solver proves optimality and returns an 18-trip plan at 59.24\,kWh, which is
0.69\% below the greedy FFD/BFD batching (59.65\,kWh) at the same trip count; the
contribution is thus a verified, physically consistent energy model rather than
a large energy margin.
\end{itemize}

% =====================================================================
\section{System Model and Problem Formulation}\label{sec:sys}
% ---------------------------------------------------------------------
\subsection{System Overview}
The temporary dispatch center $O_{01}$ provides ground control, the communication
gateway (G01, antenna 20 m above ground), and battery/energy-component
management. Three types of transportation UAVs (A/B/C) execute round trips under
a return-safety margin $\rho=20\%$, delivering box by box to 15 service areas subject to
payload, volume and energy constraints. Shadowed back-slope segments are
relayed by two relay UAVs using multi-hover-point direct transfer, maintaining
end-to-end connectivity through a ``transport--relay--G01'' two-hop link (both
directions $\ge -90$ dBm). Relay UAVs return to base to swap energy components
after each take-off/service/return cycle. Fig.~\ref{fig:system} summarizes the
overall architecture.
\begin{figure}[t]
\centering
\includegraphics[width=0.9\linewidth]{fig_system.pdf}
\caption{UAV transportation--communication collaborative system architecture.}
\label{fig:system}
\end{figure}

\textbf{Model assumptions.} Five assumptions define the scope, each trading
some realism for reproducibility. (A1)~Terrain and link state are static within
each one-second sample, matching the verification grid; shorter outages are
therefore outside the coverage guarantee. (A2)~Terrain occlusion is binary
with a fixed 10\,dB loss rather than a diffraction model, which makes the
line-of-sight test deterministic and fully reproducible but less radio-faithful.
(A3)~Wind, precipitation and equipment faults are excluded, so plans are valid
for a single deterministic post-disaster scenario. (A4)~Demand (box masses,
volumes and deadlines) is known and fixed; stochastic or rolling demand is left
to future work. (A5)~Relay demand is computed from the planned transport
trajectories, so the communication problem is conditional on the Problem-2
schedule; this preserves the stage ordering but excludes replanning after
unplanned detours.

\subsection{Physical Caliber (shared by all stages)}
\subsubsection{Distance and cruise altitude}
The great-circle distance between $O_{01}$ and service area $S_j$ uses the
Haversine formula:
\begin{equation}
\begin{aligned}
d_{Oj}&=2R\arctan\!\sqrt{\frac{a}{1-a}},\\
a&=\sin^{2}\!\frac{\Delta\varphi}{2}
+\cos\varphi_{1}\cos\varphi_{2}\sin^{2}\!\frac{\Delta\lambda}{2},
\end{aligned}
\end{equation}
where $\Delta\varphi,\Delta\lambda$ are latitude/longitude differences and $R$ is
the Earth radius. The cruise altitude along a segment takes the DEM elevation by
the nearest-pixel rule plus 50\,m:
\begin{equation}
h^{\mathrm{cruise}}_{(j)}=h^{\max}_{(j)}+50 .
\end{equation}

\subsubsection{Trip energy and return-safety margin}
The total round-trip energy consists of four parts (outbound level, return
level, outbound climb, return climb):
\begin{equation}
E^{\mathrm{total}}_{(i,j)}(q)=E^{\mathrm{go}}_{\mathrm{hor}}(q)+E^{\mathrm{ret}}_{\mathrm{hor}}(0)
+E^{\mathrm{go}}_{\mathrm{climb}}(q)+E^{\mathrm{ret}}_{\mathrm{climb}},
\end{equation}
with outbound payload $q$ and empty return. The level-flight parts follow a
nonlinear equivalent-range law: with $L_0$ the empty and $L_F$ the full-load
standard range and $Q$ the rated payload, the equivalent range at payload $q$
is
\begin{align}
L_g(q)&=L_0-(L_0-L_F)\Big(\frac{q}{Q}\Big)^{3/2},\\
E^{\mathrm{go}}_{\mathrm{hor}}(q)&=E^{\mathrm{bat}}\frac{d}{L_g(q)},\qquad
E^{\mathrm{ret}}_{\mathrm{hor}}(0)=E^{\mathrm{bat}}\frac{d}{L_0},
\end{align}
so that loading reduces range superlinearly. The climb parts are potential
energy terms with efficiency $\eta_c=0.72$ and no descent recovery:
\begin{align}
E^{\mathrm{go}}_{\mathrm{climb}}(q)&=\frac{g(m_0+q)(H-h_O)}{\eta_c\cdot3.6\times10^{6}},\\
E^{\mathrm{ret}}_{\mathrm{climb}}&=\frac{g\,m_0(H-h_S-30)}{\eta_c\cdot3.6\times10^{6}},
\end{align}
where $m_0$ is the curb mass, $H$ the cruise altitude, $h_O$ the depot
elevation, $h_S$ the service-area elevation and the $30$\,m offset is the
hand-over height above the delivery point; energies are in kWh. The usable
energy with margin $\rho$ is
\begin{equation}
E^{\mathrm{avail}}_{(i)}=E^{\mathrm{bat}}_{(i)}\cdot(1-\rho),
\qquad
E^{\mathrm{total}}_{(i,j)}(m^{\mathrm{safe}}_{(i,j)})\le E^{\mathrm{avail}}_{(i)} ,
\end{equation}
from which the maximum safe payload $m^{\mathrm{safe}}_{(i,j)}$ is obtained by
binary search. The corresponding trip time combines vertical and horizontal
flight with ground handling,
\begin{equation}
\begin{split}
t_{\mathrm{trip}}={}&(2H-h_O-h_S-30)\Big(\frac{1}{v_{\mathrm{up}}}+\frac{1}{v_{\mathrm{dn}}}\Big)
+\frac{2d}{v}\\
&+t_{\mathrm{prep}}+t^{0}_{\mathrm{hand}}
+n\,(t_{\mathrm{load}}+t^{1}_{\mathrm{hand}}),
\end{split}
\end{equation}
where $n$ is the number of boxes on the trip, $v_{\mathrm{up}},v_{\mathrm{dn}}$
are the climb and descent speeds and the last terms are the preparation,
loading and hand-over times of Table~\ref{tab:platform}.

\subsection{Notation}
Table~\ref{tab:notation} lists the main notation; unless stated otherwise, units
follow the problem data (length in m, mass in kg, time in s, energy in kWh and
power in dBm).
\begin{table}[t]
\caption{Main notation.}
\label{tab:notation}
\centering
\setlength{\tabcolsep}{4pt}
\begin{tabular}{lll}
\toprule
Symbol & Meaning & Unit \\
\midrule
$S_j$ / $O_{01}$ & service area / dispatch center & -- \\
$n_j$ & number of boxes at $S_j$ & -- \\
$\mu_k$ & UAV type of trip $k$ ($\in\{A,B,C\}$) & -- \\
$m^{\mathrm{safe}}$ & maximum safe payload & kg \\
$\rho$ & return-safety margin & -- \\
$N_{\mathrm{trips}}$ & number of round trips & -- \\
$T_{\mathrm{total}}$ & cumulative operating time & s \\
\midrule
$E_{\mathrm{hor}},E_{\mathrm{climb}},E_{\mathrm{total}}$ & level/climb/total energy & kWh \\
$L_g(q)$ & nonlinear equivalent range & m \\
$L_0$, $L_F$ & empty / full-load standard range & m \\
$m_0$, $Q$ & curb mass / rated payload & kg \\
$v_{\mathrm{up}}$, $v_{\mathrm{dn}}$ & climb / descent speed & m/s \\
$\eta_c$ & climb energy efficiency ($0.72$) & -- \\
$h^{\mathrm{cruise}}$ & cruise altitude (max elev.\ $+50$\,m) & m \\
SOC & state of charge (return $\ge20\%$) & -- \\
$C_{\max}$ & mission completion time (latest return) & s \\
$L_{\mathrm{soft}}$ & accumulated soft delay & min \\
RSRP & received power (threshold $-90$\,dBm) & dBm \\
$\mathcal{P}=\{P_1,\dots,P_K\}$ & task partition ($K{=}2$ or $3$) & -- \\
$N(P_i,q)$ & peak concurrent demand of $q$ & units \\
$D$, $G$, $J$ & scale/gap/balance & mixed \\
$I_q$ & resource inventory vector & units \\
\bottomrule
\end{tabular}
\end{table}

% =====================================================================
\section{Methodology}\label{sec:method}
% ---------------------------------------------------------------------
\subsection{Overview of the Unified Framework}
We establish a ``transport--communication--resource'' collaborative
optimization chain: with 30\,m DEM terrain, 15 service areas, 80 indivisible
boxes and heterogeneous UAV fleets plus link parameters as unified inputs, we
successively build DEM-ELBPP (Problem 1), HF-VRPBT (Problem 2), TD-CJO (Problem
3) and SC-MAE-RG (Problem 4). The four layers introduce return-safety margin and
batching, deadline concurrency and battery turnover, full-process continuous
communication, and partitioning resource allocation, respectively; the overall
technical route is shown in Fig.~\ref{fig:route}.
\begin{figure}[t]
\centering
\includegraphics[width=0.9\linewidth]{fig_route.pdf}
\caption{Overall technical route of the proposed framework.}
\label{fig:route}
\end{figure}

\subsection{Problem 1: DEM-Driven Energy-Aware Bin Packing (DEM-ELBPP)}
For the box set $B_j$ at service area $S_j$ ($|B_j|=n_j$), batching is modeled as
an energy-constrained hierarchical bin packing: each trip $\pi_k$ loads a nonempty
subset of $B_j$ executed by type $\mu_k\in\{A,B,C\}$, subject to payload, volume
and energy constraints. Among the three objectives (trip count, total energy,
cumulative operating time) we adopt \emph{lexicographic} optimization instead of a
weighted sum, because the three have different dimensions and arbitrary weights
would change the optimum subjectively:
\begin{equation}
\min_{\text{lex}}\Big(N_{\mathrm{trips}},\ \sum_{k}E_{k},\ T_{\mathrm{total}}\Big)
\end{equation}
subject to
\begin{align}
&\bigcup_{k}\pi_k=B_j,\qquad \pi_k\cap\pi_{k'}=\varnothing\ (k\ne k'),\\
&\sum_{n\in\pi_k}m_{n}\le m^{\mathrm{safe}}_{(i,j)},\qquad
\sum_{n\in\pi_k}V_{n}\le V^{\max}_{(i)},\\
&\mu_k\in\{A,B,C\},\quad \pi_k\subseteq B_j\ \text{(single area, no crossing).}
\end{align}
The key idea is to first back-solve the safe payload from the energy equation and
then perform set partitioning under the payload--volume--energy triple constraint,
discretizing the continuous energy constraint into bin-packing feasibility.

Candidate batches are enumerated per area (the largest area holds $n_j=15$
boxes, so at most $2^{15}-1=32767$ subsets), and the exact solution uses
set-partitioning dynamic programming:
\begin{equation}
dp[\text{mask}\cup\pi]=\min\Big(dp[\text{mask}\cup\pi],\ dp[\text{mask}]\oplus(\pi)\Big),
\end{equation}
where $\oplus$ is vector addition and triples are compared lexicographically;
the optimal batch sequence is then recovered by backtracking the predecessor
chain. Two implementation facts bound the computational cost. First, the 45
machine--area safe payloads are precomputed once by a 60-iteration binary
search on the energy equation, so the DP never evaluates the continuous energy
model. Second, because each area holds at most fifteen boxes, the per-area
candidate count is at most $2^{15}-1=32767$ subsets before the
payload--volume--energy filter and far fewer after it, and the DP state space
is the $2^{15}$ masks of that area; mixed-type batching within an area is
allowed, so the DP chooses the cheapest machine per candidate subset rather
than fixing the type in advance. The overall procedure is summarized in Algorithm~\ref{alg:q1}.

\begin{algorithm}[t]
\caption{Exact batching (subset enumeration + set-partitioning DP)}
\label{alg:q1}
\begin{algorithmic}[1]
\STATE \textbf{input} box sets $B_j$, UAV parameters, DEM, margin $\rho$
\STATE binary-search $m^{\mathrm{safe}}_{(i,j)}$ for each machine/area
\STATE enumerate feasible subsets per area, filter by payload/volume/energy $\to$ candidate batches $C_j$
\STATE init $dp[0]=(0,0,0)$, others $\infty$; relax candidate batches lexicographically
\STATE backtrack to output the optimal batching plan
\STATE scan $\rho=20\%\sim40\%$ and repeat for sensitivity
\STATE \textbf{output} 18-trip plan, per-area machine/energy, margin sensitivity
\end{algorithmic}
\end{algorithm}

\subsection{Problem 2: Heterogeneous Fleet Routing with Battery and Time Windows (HF-VRPBT)}
Problem 2 relaxes the single-area structure and adds resource and time-window
constraints: each trip may visit one or more service areas before returning. The
decisions jointly determine box batching, visit sequence, machine type
$\mu_k$, the executing airframe, the shared battery and the trip start time
$t_0$, and a trip occupies the whole interval from preparation through
take-off, flight, delivery, return and charging. With $B$ the 80 boxes,
$B^{\mathrm{hard}}$ the 31 hard-deadline boxes and $B^{\mathrm{soft}}$ the rest,
the constraints enforce one delivery per box, per-trip payload and volume
limits, segment-wise residual payload, trip energy, return SOC $\ge\rho$, hard
deadlines $t_b\le d_b$ for all $b\in B^{\mathrm{hard}}$, airframe and battery
mutual exclusion, and recharge to $100\%$ before battery reuse. Formally, with
$x_{kb}\in\{0,1\}$ marking the assignment of box $b$ to trip $k$ and
$[s_k,e_k]$ the interval that trip $k$ occupies its airframe and battery, the
scheduler accepts a plan only if
\begin{equation}
\begin{aligned}
&\textstyle\sum_{k}x_{kb}=1\ \ (b\in B),\qquad
\sum_{b}x_{kb}m_{b}\le m^{\mathrm{safe}}_{(\mu_k)},\\
&\textstyle\sum_{b}x_{kb}V_{b}\le V^{\max}_{(\mu_k)},\qquad
t_{b}\le d_{b}\ \ (b\in B^{\mathrm{hard}}),\\
&\mathrm{SOC}_{k}=1-E_{k}/E^{\mathrm{bat}}_{(\mu_k)}\ge\rho,\\
&[s_k,e_k]\cap[s_{k'},e_{k'}]=\varnothing
\ \ \text{(shared airframe or battery),}
\end{aligned}
\end{equation}
i.e.\ the occupied intervals of trips sharing an airframe or a battery do not
overlap, and a battery becomes reusable only after a full recharge. Four
objectives---timeliness, makespan $C_{\max}$ (latest return), transport energy
and trip count---are compared in a non-dominated sense, with the soft delay
$L_{\mathrm{soft}}=\sum_{b\in B^{\mathrm{soft}}}\max(0,t_b-d_b)$ measuring the
lateness of the remaining boxes. Since the full trip-routing space is
intractable, we enumerate a restricted candidate set (whole-area trips inherited
from Problem 1, split first-batch tasks and short cross-area tours) and apply a
deterministic urgency-first greedy schedule that inserts each trip for the task
whose hard deadline is nearest while updating the affected airframe and battery
occupancies. The three candidate classes play distinct roles: the inherited
whole-area trips preserve the energy-optimality of the Problem-1 batches, the
split first-batch tasks decouple the urgent first deliveries from the residual
load of the same area so that hard deadlines can be met before the full batch
would be ready, and the short cross-area tours absorb small residual loads that
would otherwise each cost a dedicated round trip. At every insertion the
scheduler recomputes the occupancy intervals of the chosen airframe and its
battery, rejects placements that violate the mutual-exclusion or recharge
conditions of the constraint set above, and leaves soft-deadline lateness to be
priced by $L_{\mathrm{soft}}$ rather than blocking on it. The search evaluates 57
candidate plans and yields five non-dominated plans (Table~\ref{tab:pareto}):
19 trips ($66.30$\,kWh), the recommended 20-trip plan ($65.25$\,kWh,
$C_{\max}=2.72$\,h, soft delay $137.6$\,min), two further 20-trip plans that
trade a longer makespan for smaller soft delay, and 21 trips ($66.00$\,kWh);
the recommended 20-trip plan
dominates on energy and completion time, uses all eight airframes and satisfies
all 31 hard-deadline boxes.

\subsection{Problem 3: Task-Driven Joint Transportation--Communication Scheduling (TD-CJO)}
\label{sec:mhprs}
\subsubsection{Link feasibility}
Link feasibility uses a free-space path-loss model with a binary DEM
line-of-sight occlusion test (10\,dB occlusion loss when blocked):
\begin{equation}
\begin{aligned}
\mathrm{RSRP}=\ &P_t+G_t+G_r-32.45-20\log_{10}f-20\log_{10}d\\
&-L_{\mathrm{obs}}-L_{\mathrm{sys}},
\end{aligned}
\end{equation}
with $L_{\mathrm{obs}}\in\{0,10\}$\,dB and $L_{\mathrm{sys}}=3$\,dB. Both
directions (transport$\leftrightarrow$relay and relay$\leftrightarrow$G01) must
satisfy $\mathrm{RSRP}\ge P_{\mathrm{thr}}=-90$\,dBm. The relay platform is a
$23.5$\,kg multi-rotor with a $3.2$\,kWh swappable energy component: it cruises
at $1.15$\,kW, hovers at $1.05+0.05$\,kW per active served link, recharges in
$1800$\,s under the two-stage charging rule, and must land with
$\mathrm{SOC}\ge0.20$; six components circulate between the two airframes and
the chargers.

\subsubsection{Multi-hover-point relay scheduling (MHPRS)}
\textbf{Definition (MHPRS).} Given a set of relay-required time intervals and a
set of candidate hover points, a \emph{relay sortie} is a take-off/service/return
cycle of one relay UAV serving an ordered sequence of hover points connected by
direct-transfer flights; MHPRS is to cover all relay-required points by the
minimum-cost set of sorties under link-budget, hover-height, resource-exclusion,
return-SOC and component-inventory constraints.
Classical relay placement asks where to install fixed relays so that a static
set of terminals becomes connected, and it optimizes over positions only. MHPRS
differs in three respects: (i) the demand is \emph{time-varying}, because
relay-required points appear and disappear as the transport trips progress, so
coverage is a spatio-temporal rather than a purely spatial requirement; (ii) a
single relay airframe may serve an \emph{ordered sequence} of hover points
within one sortie, and the transfer flight and link re-establishment between
consecutive points consume time and energy that must enter the constraints; and
(iii) energy, hover-height and component-turnover limits bind the solution,
since only two relay airframes share a finite inventory of six energy
components. Formally, the constraints are the full-process coverage condition
(each transport point is either directly connected to G01 or lies in some
sortie's service interval with both hops above threshold), relay-before-task
$s_c\le\min\{t_p\}$, bidirectional link budget
$\mathrm{RSRP}\ge P_{\mathrm{thr}}$, hover-height limit
$h_c\le h^{\mathrm{terr}}_c+300$\,m, resource exclusion, return SOC $\ge0.20$,
component inventory (in-flight plus charging components $\le6$), coverage rate
$\ge\eta_{\mathrm{target}}$, and the transfer--link-establish timing between
consecutive segments. The coverage rate is the fraction of verified
relay-required trajectory points, counted after removing duplicate
boundary samples between consecutive flight segments so that no point is
counted twice.

The solution procedure mirrors this structure. Relay-required points are
clustered into service intervals by temporal gaps larger than $400$\,s;
candidate hover points are generated on the $0.01^\circ$ DEM grid within
$300$\,m above ground and $12$\,km of the depot, keeping only points from
which every relay-required point of a segment passes the bidirectional link
budget. A greedy pass selects the fewest candidates per interval, splitting
the interval at its median time whenever a segment remains uncovered, and the
resulting adjacent segments are packed into sorties under the component
inventory and return-SOC constraints with earliest-available assignment of
the two airframes.

\textbf{Proposition 1 (mode dominance).} Let $\mathcal{S}_{\mathrm{SH}}$ be
the set of coverage plans feasible when every sortie is restricted to a single
hover point, and $\mathcal{S}_{\mathrm{MH}}$ the set feasible under MHPRS with
ordered hover-point chains. Then
$\mathcal{S}_{\mathrm{SH}}\subseteq\mathcal{S}_{\mathrm{MH}}$; consequently,
for any coverage target $\eta$, (i) if $\eta$ is attainable at all by the
single-hover mode with a given fleet, it is attainable by MHPRS with the same
fleet, and (ii) the minimum energy required under MHPRS is no greater than
under the single-hover mode. \emph{Proof.} A single-hover sortie is the
length-one special case of an MHPRS chain: relabeling every single-hover
sortie as a length-one chain adds no transfer flight, and all other
constraints (link budget, hover height, resource exclusion, return SOC,
component inventory) are identical, so the relabeled plan remains feasible at
the same energy. Minimizing a common cost over a superset of feasible plans
cannot raise the optimum.\hfill $\square$

\noindent
\emph{Empirical gap.} The proposition is a feasibility statement; how much the
enlarged feasible set buys on the test instance is measured empirically. The
single-hover mode covers only $51.0\%$, $62.7\%$ and $82.1\%$ of the
relay-required points with one, two and three airframes, so the single-hover
mode does not attain full coverage for $N\le3$; under MHPRS the same two
available airframes chain several hover points per sortie and the verified plan
covers all 11683 sampled relay-required points with only four sorties, showing
where the practical gap comes from.

The combined stage-3 objective is
\begin{align}
\min Z &= E_{\mathrm{trans}}+E_{\mathrm{relay}}, \notag\\
E_{\mathrm{relay}} &= \sum_{c}\!\Big(E^{\mathrm{fly}}_{c}
+(P_{\mathrm{hover}}+P_{\mathrm{comm}})\tfrac{e_c-s_c}{3600}
+P_{\mathrm{hover}}\tfrac{w_c}{3600}\Big),
\end{align}
where $w_c$ is the accumulated hovering wait time within sortie $c$ and
$E^{\mathrm{fly}}_{c}$ is the outbound/transfer/return flight energy, subject to
the constraint set listed in the MHPRS definition above. The resulting
multi-hover-point direct-transfer mode, together with its coverage advantage
over the single-hover baseline, is illustrated in Fig.~\ref{fig:mhprs} and
summarized in Algorithm~\ref{alg:q3}.

\begin{algorithm}[t]
\caption{Task-driven relay assignment (TD-CJO)}
\label{alg:q3}
\begin{algorithmic}[1]
\STATE \textbf{input} Problem-2 schedule, DEM, link params, 2 relay UAVs, 6 components
\STATE sample trajectories per second, mark relay-required points
\STATE cluster relay-required points by time ($>$400\,s gap) into service intervals
\STATE generate candidate hover points (DEM-covered, $\le$300\,m, $\le$12\,km)
\STATE greedy select points per cluster; split by median time if uncovered
\STATE pack adjacent segments into sorties (multi-hover direct transfer); earliest-available assignment
\STATE check component turnover (two-stage charging) and return SOC
\STATE full verification by per-second sampling and boundary checks
\STATE \textbf{output} 2-UAV / 4-sortie main plan (11683 sampled relay points, 100\% covered)
\end{algorithmic}
\end{algorithm}

\begin{figure}[t]
\centering
\includegraphics[width=0.9\linewidth]{fig_mhprs.pdf}
\caption{Multi-hover-point direct-transfer relay scheduling mode.
(a) Hover-point chains of the two relay sorties, where the direct transfer
between hover points is the decisive factor; (b) coverage of relay-required
points versus the number of relay airframes under the single-hover round-trip
baseline, compared with the proposed MHPRS.}
\label{fig:mhprs}
\end{figure}

\subsection{Problem 4: Partition-Based Resource Accounting (SC-MAE-RG)}
Problem 4 takes the Problem-3 joint plan as input and groups the 15 service
areas into $K\in\{2,3\}$ task groups, subject to each area belonging to exactly
one group, each group containing at least one area, and any areas that share a
transport trip being placed in the same group. Each group independently accounts
for the eight resource classes (A/B/C airframes, A/B/C batteries, relay airframes
and energy components) at their peak concurrency, because resources may not be
reassigned across groups. For group $P_i$ and resource class $q$ the peak
concurrent demand is
$N(P_i,q)=\max_\tau\sum_{u:S_u\cap P_i\ne\varnothing}\mathbf{1}\{\tau\in[a_u,b_u],\,\mathrm{type}(u)=q\}$,
from which the configuration scale $D=\sum_i\sum_q N(P_i,q)$, the inventory gap
$G=\sum_q\max(0,\sum_i N(P_i,q)-I_q)$, the cross-group relay count, the
replication total and the inter-group balance index $J$ follow; candidate
partitions are then ranked lexicographically by the key (unique cross-group
relays, extra cross-group replicas, $D$, $G$, $J$), which penalizes relay
replication before scale. The balance index is the mean
coefficient of variation across six workload dimensions (boxes, mass, sorties,
energy, busy duration and relay count),
\begin{equation}
J(\mathcal{P})=\frac{1}{6}\sum_{\ell=1}^{6}\frac{\sigma_{\ell}(\mathcal{P})}
{\mu_{\ell}(\mathcal{P})},
\end{equation}
so that $J{=}0$ marks a perfectly even split and larger values indicate
workload concentration. The $K{=}2$ partition $[11,4]$
gives $D{=}31$ and $G{=}5$ (two B-type airframes, two B-type batteries and one
relay airframe), whereas $K{=}3$ $[7,4,4]$ gives $D{=}38$ and $G{=}10$; the
balance index is $J{=}0.613$ for $K{=}2$ and $J{=}0.438$ for $K{=}3$, so $K{=}2$
is preferred as it is smaller and less redundant. A zero-replication control
lowers the scale to $D{=}30$ ($K{=}2$) and $D{=}31$ ($K{=}3$), showing that
cross-group relay replication is the main source of redundancy. An
entropy-weighted TOPSIS ranking identifies S001 as the most vulnerable service
area, consistent with the heavy load carried by the larger group.

% =====================================================================
\section{Experimental Design and Results}\label{sec:exp}
% ---------------------------------------------------------------------
\subsection{Setup}
All results are computed on the standard mountainous test instance supplied with
the problem: one dispatch center $O_{01}$ (with gateway G01), 15 service areas
$S_1$--$S_{15}$, 80 indivisible boxes of four supply categories, three transport
UAV types with a shared battery pool, two relay UAVs with six energy components,
a 30\,m DEM and the specified wireless-link parameters. For the batching stage
the exact set-partitioning solver is compared with the greedy FFD and BFD
heuristics over $N{=}30$ randomized instances derived from the real case; for the routing and relay stages the
multi-hover-point plan is compared with single-area, single-hover-point and
augmented-fleet controls. Metrics are the number of trips, total energy,
makespan, soft delay and communication coverage; all solvers use Python
(NumPy/pandas) and the reported values are reproduced by an independent
verification script. Relay-required points are counted on per-second sampled
trajectories (1\,s step), and the coverage figures quoted here for the main plan
refer to this caliber, whereas the threshold and clearance scans reported below
resample the same trajectories at a coarser 5\,s search step, which is why their
relay-required counts are smaller. The fleet and energy parameters are listed in
Table~\ref{tab:platform} and the communication and terrain parameters in
Table~\ref{tab:link}; the return-SOC floor is $\rho=20\%$ and the climb energy
efficiency is $0.72$. The 80 boxes span four categories: medical kits
($3$\,kg, $0.012$\,m$^3$), drinking water ($14$\,kg, $0.027$\,m$^3$),
emergency food ($8$\,kg, $0.028$\,m$^3$) and hygiene items ($6$\,kg,
$0.035$\,m$^3$).

\begin{table}[t]
\centering
\footnotesize
\setlength{\tabcolsep}{3pt}
\caption{Heterogeneous fleet and energy parameters of the test instance.}
\label{tab:platform}
\begin{tabular}{lccc}
\toprule
Parameter & Type A & Type B & Type C \\
\midrule
Payload $Q$ (kg) & 25 & 30 & 80 \\
Curb mass $m_0$ (kg) & 70.0 & 65.0 & 69.9 \\
Cabin volume (m$^3$) & 0.06 & 0.073 & 0.25 \\
Cruise speed $v$ (m/s) & 12 & 15 & 15 \\
Climb/descent speed $v_{\mathrm{up}}/v_{\mathrm{dn}}$ (m/s) & 3.0/2.5 & 3.0/2.5 & 2.5/2.0 \\
Battery capacity $E^{\mathrm{bat}}$ (kWh) & 4.5 & 4.0 & 8.0 \\
Range empty/full $L_0/L_F$ (km) & 25/20 & 28/16 & 26/12 \\
Airframes & 4 & 2 & 2 \\
Shared batteries & 6 & 4 & 4 \\
Full-charge time (s) & 1800 & 2400 & 3000 \\
Preparation $t_{\mathrm{prep}}$ (s) & 300 & 300 & 300 \\
Loading per box $t_{\mathrm{load}}$ (s) & 30 & 30 & 30 \\
Hand-over base${}+{}$per box (s) & $150{+}30$ & $150{+}30$ & $180{+}36$ \\
\bottomrule
\end{tabular}
\end{table}

\begin{table}[t]
\centering
\footnotesize
\setlength{\tabcolsep}{3pt}
\caption{Communication and terrain parameters of the test instance.}
\label{tab:link}
\begin{tabular}{ll}
\toprule
Parameter & Value \\
\midrule
Carrier frequency & 2400\,MHz \\
DEM resolution & 30\,m \\
Occlusion loss (blocked line of sight) & 10\,dB \\
System loss $L_{\mathrm{sys}}$ & 3\,dB \\
Receiver sensitivity / design margin & $-98$\,dBm / 8\,dB \\
Receiving threshold $P_{\mathrm{thr}}$ & $-90$\,dBm \\
Transport UAV transmit power / antenna gain & 20\,dBm / 3\,dBi \\
Relay access link (to transport UAVs) & 20\,dBm / 6\,dBi \\
Relay backhaul link (to G01) & 19\,dBm / 8\,dBi \\
Gateway G01 (20\,m mast) & 27\,dBm / 12\,dBi \\
\bottomrule
\end{tabular}
\end{table}

\subsection{Baseline Comparison}
Table~\ref{tab:q1_baseline} compares the exact dynamic-programming batching with
the classical FFD, BFD and WFD heuristics over $N{=}30$ randomized instances
derived from the real case. On average the proposed method returns the fewest
trips and the lowest energy (67.18\,kWh), and the $2.49\%$ mean energy penalty of
the heuristics is statistically significant ($p<10^{-4}$, two-sided Wilcoxon
signed-rank test). On the real instance the same ordering holds with a much
smaller gap (59.24\,kWh against 59.65\,kWh, i.e.\ $0.69\%$): the advantage of the
exact solver lies in consistent behaviour across instances rather than in a
single large energy margin. The randomized instances are generated with the
fixed seed 20251015: the box count is drawn uniformly from
$\{40,\dots,120\}$, the service-area coordinates are jittered with
$\sigma=0.0005^\circ$, and the box categories are resampled from the empirical
distribution of the real case. The Wilcoxon statistic uses the 29 non-zero
per-instance pairs, giving an exact two-sided $p=2.56\times10^{-6}$ for both
FFD and BFD. Across the family the per-instance energy spans $42.6$--$89.3$\,kWh
and the trip count $14$--$22$, so the comparison covers instances both smaller
and larger than the real case rather than a narrow neighborhood of it.

\begin{table}[t]
\centering
\footnotesize
\setlength{\tabcolsep}{3pt}
\caption{Comparison of the proposed exact dynamic-programming grouping against classical bin-packing heuristics over $N=30$ randomized instances derived from the real case. Sorties and total energy (in kWh) are reported as mean$\pm$std; the improvement is the mean energy increase relative to the proposed method; $p$-values are from a two-sided Wilcoxon signed-rank test on per-instance total energy (proposed vs.\ each heuristic).}
\label{tab:q1_baseline}
\begin{tabular}{lcccc}
\toprule
Method & Sorties & Energy & Improv. & $p$-value \\
\midrule
Proposed (Exact DP) & 18.5$\pm$2.3 & 67.18$\pm$14.32 & --- & --- \\
FFD & 18.6$\pm$2.4 & 68.85$\pm$15.26 & $+2.49\%$ & $<10^{-4}$ \\
BFD & 18.6$\pm$2.4 & 68.85$\pm$15.26 & $+2.49\%$ & $<10^{-4}$ \\
WFD (worst-fit) & 18.5$\pm$2.3 & 69.06$\pm$15.23 & $+2.79\%$ & $<10^{-4}$ \\
\bottomrule
\end{tabular}
\end{table}

\subsection{Ablation and Sensitivity}
The models are examined by ablation and one-factor sensitivity analysis.

For the routing stage, the candidate search evaluates 57 plans, of which five
are non-dominated (Table~\ref{tab:pareto}). The recommended 20-trip main plan
has the lowest makespan ($2.72$\,h) and near-minimal energy ($65.25$\,kWh);
its soft delay of $137.6$\,min concerns only soft-deadline boxes, because all
31 hard-deadline boxes are delivered on time. Decomposed by aircraft type, the
main plan flies four A-type trips ($9.59$\,kWh), seven B-type trips
($16.32$\,kWh) and nine C-type trips ($39.34$\,kWh), so the heavy C type
carries $60.3\%$ of the transport energy. All eight airframes are used: the
four A airframes fly one trip each, while the two B and two C airframes fly up
to four and five trips respectively under battery-turnover constraints.
Nineteen trips serve a single service area and only one trip (T18) chains two
areas (S011 then S008), which shows that cross-area chaining is profitable only
where the inherited Problem-1 batches leave small residual loads. The makespan
is set by the last C-type trip to S007, which returns at $9794$\,s.

To measure the optimality gap of the scheduling layer independently of the
candidate-structure search, the urgency-first scheduler and an exact CP-SAT
model (OR-Tools) are given the \emph{same} fixed 20-trip task set and the
\emph{same} per-trip feasible-type choice set; CP-SAT jointly selects
aircraft type, numbered airframe, battery and start time under no-overlap and
two-stage recharge constraints and lexicographically minimizes hard-deadline
violations, $C_{\max}$, energy and soft delay. Durations are rounded \emph{up}
to whole seconds, so each CP-feasible schedule is strictly feasible under the
continuous physics, and the returned schedule is additionally re-checked
against the continuous energy/time model (zero violations).
Table~\ref{tab:benchq2}: CP-SAT proves optimality in $0.53$\,s, shortening
$C_{\max}$ from $2.46$ to $2.11$\,h ($-14.3\%$), energy by $4.6\%$ and soft
delay from $126.5$ to $5.3$\,min; the task set and physical options are
identical, so the gap is attributable to the scheduler. On 30 synthetic
instances drawn from the empirical duration/energy/deadline distributions
(infeasible ones resampled), CP-SAT wins all 30, $2.47{\pm}0.33$ to
$2.07{\pm}0.28$\,h (median gain $15.2\%$, Wilcoxon $p{<}10^{-8}$), at $5.0$\,s
mean wall time, proving optimality on 28/30 within 30\,s. CP-SAT solves only
the scheduling sub-problem: it receives batched trips and precomputed physical
options, leaving DEM feasibility and batching-structure enumeration to our chain.

\begin{table}[t]
\centering
\footnotesize
\setlength{\tabcolsep}{4pt}
\caption{External solver benchmark for the Problem-2 scheduling layer:
20 fixed trips, identical aircraft-type choice set for both schedulers
(wall time, 8-core CPU). CP-SAT was proven optimal; its schedule independently
re-verified against the continuous model with zero constraint violation.}
\label{tab:benchq2}
\begin{tabular}{lcccc}
\toprule
Scheduler & $C_{\max}$ (h) & Energy (kWh) & Soft (min) & Wall time \\
\midrule
Urgency-first greedy & 2.46 & 67.38 & 126.5 & $<$0.01\,s \\
CP-SAT (OR-Tools)    & \textbf{2.11} & \textbf{64.31} & \textbf{5.3} & 0.53\,s \\
\bottomrule
\end{tabular}
\end{table}

\begin{table}[t]
\centering
\footnotesize
\setlength{\tabcolsep}{3pt}
\caption{Non-dominated routing plans (five out of 57 evaluated candidates).
The recommended main plan is marked in bold.}
\label{tab:pareto}
\begin{tabular}{cccc}
\toprule
Trips & Energy (kWh) & Makespan (h) & Soft delay (min) \\
\midrule
19 & 66.30 & 2.94 & 105.77 \\
\textbf{20} & \textbf{65.25} & \textbf{2.72} & \textbf{137.61} \\
20 & 65.57 & 2.80 & 109.95 \\
20 & 67.05 & 2.80 & 97.56 \\
21 & 66.00 & 2.80 & 95.50 \\
\bottomrule
\end{tabular}
\end{table}

For the batching model, forcing a single aircraft type raises the energy to
112.39\,kWh (all~A), 68.48\,kWh (all~B) or 74.06\,kWh (all~C), against
59.24\,kWh for the mixed exact plan, so the aircraft--payload--range match is
the dominant factor. In the exact plan twelve areas are served by a single
trip and three areas need two (S001--S003); the binding energy constraint is
the C-type trip to S004, which returns with SOC $0.228$, only $2.8$ percentage
points above the margin floor, while the DEM-driven cruise altitudes of the 18
trips range from $273$\,m (S001) to $592$\,m (S014).
Table~\ref{tab:margin} scans the return-safety margin: the
18-trip, 59.24\,kWh plan is unchanged for $\rho\le20\%$, and feasibility breaks
at $\rho{=}40\%$, where the empty round trips to S004 and S008 exceed the
equivalent range.

\begin{table}[t]
\centering
\footnotesize
\setlength{\tabcolsep}{3pt}
\caption{Return-safety margin scan for the batching stage (Problem 1).}
\label{tab:margin}
\begin{tabular}{ccccl}
\toprule
$\rho$ & Trips & Energy (kWh) & Duration (h) & Remark \\
\midrule
10\%--20\% & 18 & 59.24 & 9.11 & recommended $\rho{=}20\%$ \\
25\% & 19 & 61.19 & 9.61 & \\
30\% & 20 & 67.36 & 10.17 & \\
35\% & 25 & 78.19 & 12.65 & \\
40\% & --- & --- & --- & infeasible (S004, S008) \\
\bottomrule
\end{tabular}
\end{table}

For the relay stage, Table~\ref{tab:threshold} tightens the receiving
threshold: from $-98$ to $-95$ to $-90$\,dBm the relay-required points grow
from 0 to 806 to 2311 and the directly served areas shrink from 15 to 10 to 3.
Under the single-hover-point round-trip baseline only 51.0\%, 62.7\% and
82.1\% of these points are covered by one, two and three airframes, confirming
that the transfer between hover points is essential. Raising the safety
clearance above the terrain-following cruise altitude from 0 to 25 to 50\,m
increases the direct-connection rate only moderately, from 61.7\% to 63.2\% to
64.7\%, so the relay requirement is governed mainly by the receiving threshold
rather than by the clearance margin. The independently rebuilt four-sortie plan
assigns 5137, 508, 1558 and 4480 relay-required points to the four sorties,
with a rebuilt relay energy of $5.52$\,kWh against the frozen $5.43$\,kWh (a
$1.6\%$ deviation), and the weakest verified link measures $-89.66$\,dBm,
only $0.34$\,dB above the $-90$\,dBm threshold. Table~\ref{tab:sortie} breaks the
verification down per sortie: the two multi-hover sorties (R01M1 with three
hover points, R02M2 with two) carry $82.3\%$ of the relay-required points and
show the largest, still small, rebuild deviations, because their inter-hover
transfer segments had to be reconstructed independently --- the frozen schedule
registers only one hover coordinate per sortie --- whereas the single-hover
sorties match the frozen energy within $0.2\%$--$2.4\%$. The reconstruction
rule behind these rebuilt values is conservative: the frozen point log contains
11683 rows with 68 boundary duplicates (11615 unique sortie--time pairs), and
the two missing transfer hover points were re-searched on the $0.01^\circ$ DEM
grid within $300$\,m above ground and $12$\,km of the depot, taking among all
candidates that cover the whole segment the one nearest to the segment's
point-cloud centroid. Neither hover geometry nor relay energy was reverse-
engineered from the frozen plan; only the transfer-climb efficiency was
calibrated ($\eta_c\approx0.50$ vs.\ the design value $0.72$), so the frozen plan
is unaffected.

\begin{table}[t]
\centering
\footnotesize
\setlength{\tabcolsep}{3pt}
\caption{Per-sortie verification of the four-sortie relay plan. Points are logged
one-second samples (11615 unique after deduplication); energy is the rebuilt value
against the frozen log (kWh); SOC is the rebuilt return state of charge; R01 flies
M1--M2, R02 flies M1--M2.}
\label{tab:sortie}
\begin{tabular}{llcccc}
\toprule
Sortie & Hover-point chain & Points & Energy (kWh) & SOC \\
\midrule
R01M1 & PER1\_E$\to$PER2$\to$PER3\_NE & 5137 & 1.844 / 1.823 & 0.424 \\
R01M2 & single hover point & 508 & 1.081 / 1.056 & 0.662 \\
R02M1 & single hover point & 1558 & 0.703 / 0.701 & 0.781 \\
R02M2 & PER3\_SW$\to$PER4 & 4480 & 1.892 / 1.853 & 0.409 \\
\midrule
Total & 7 hover points, 7 service windows & 11683 & 5.519 / 5.433 & -- \\
\bottomrule
\end{tabular}
\end{table}

\begin{table}[t]
\centering
\footnotesize
\setlength{\tabcolsep}{3pt}
\caption{Receiving-threshold scan for the relay stage (5\,s search sampling,
6028 sampled trajectory points).}
\label{tab:threshold}
\begin{tabular}{lccc}
\toprule
Threshold & Relay-required & Direct-connected & Direct-served areas \\
\midrule
$-90$\,dBm & 2311 & 3717 & 3 \\
$-95$\,dBm & 806 & 5222 & 10 \\
$-98$\,dBm & 0 & 6028 & 15 \\
\bottomrule
\end{tabular}
\end{table}

For the partitioning stage, the $K{=}2$ plan groups S001--S009, S011 and S015
(11 areas, 68 boxes, 15 sorties, $54.37$\,kWh) against S010 and S012--S014 (4
areas, 12 boxes, 5 sorties, $10.87$\,kWh), with scale $D{=}31$, gap $G{=}5$ and
balance index $J{=}0.613$; the $K{=}3$ plan splits $9/6/5$ sorties
($35.87/18.50/10.87$\,kWh) with $D{=}38$, $G{=}10$ and $J{=}0.438$. The
zero-replication control lowers the scale to $D{=}30$ and $D{=}31$
respectively, isolating cross-group relay replication as the cause of the
extra redundancy. The entropy-weighted TOPSIS vulnerability ranking behind the
S001 identification uses criteria weights of $0.412$ (mass), $0.385$ (volume),
$0.106$ (distance) and $0.097$ (elevation); the five most vulnerable areas are
listed in Table~\ref{tab:topsis}, where S001 ranks first by a factor of about
$1.9$ in demand because it combines the largest demand with the shortest distance,
making it the natural first wave. Table~\ref{tab:gap} details the
peak concurrent demand of the recommended $K{=}2$ partition: the smaller group
still requires its own B-type fleet because trips may not be reassigned across
groups, and exactly one relay sortie (R01M1) serves areas of both groups, which
forces the third relay airframe and explains the replication term. The
lexicographic key confirms this: $K{=}2$ needs one unique cross-group relay
sortie and one extra replica, against three of each for $K{=}3$, and no
evaluated candidate achieves fewer cross-group relays for either $K$, so the
comparison is decided on scale and gap. The
single-hover baseline of Fig.~\ref{fig:mhprs}(b) is evaluated under exactly
the same link budget and airframe inventory as MHPRS, so the coverage gap is
attributable to the mode itself rather than to more generous resources.

The resource-directed enumeration is benchmarked against generic unsupervised
clustering: $K$-means ($100$ restarts) and Ward agglomerative clustering on
standardized coordinates (geo) and on coordinates plus
boxes--mass--volume--first-batch demand (geo+demand). Each clustering output
receives the identical S008--S011 feasibility repair (shared sortie T18) and
is scored by the identical peak-concurrency accounting. Among operationally
valid partitions (every group $\ge$3 areas), Table~\ref{tab:benchq4} shows
SC-MAE-RG attains the smallest scale and gap for both $K$: $31/5$ for $K{=}2$
against $36/9$ and $37/9$, and $38/10$ for $K{=}3$ against $43/15$ and
$40/12$. Clustering obtains smaller (better) $J$ but at the price of more
cross-group relay sorties and replicas. With demand features, both methods
isolate S001 (154\,kg, $1.9\times$ the next-largest area) into a
singleton group ($1{+}14$ for $K{=}2$; $1{+}7{+}7$ and $1{+}5{+}9$ for
$K{=}3$), violating the support-unit rule, so their apparently lower $D/G$
is inoperable: generic clustering minimizes feature dispersion without
perceiving peak concurrency or the non-splittable shared-sortie coupling.

\begin{table}[t]
\centering
\footnotesize
\setlength{\tabcolsep}{4pt}
\caption{External clustering benchmark for Problem 4 under identical resource
accounting ($D$: total peak demand; $G$: inventory gap; $J$: balance index,
smaller is more balanced; cross: cross-group relay sorties).}
\label{tab:benchq4}
\begin{tabular}{llcccc}
\toprule
Method & $K$ & $D$ & $G$ & $J$ & Cross \\
\midrule
$K$-means (geo)               & 2 & 36 & 9  & 0.17 & 2 \\
Ward (geo)                    & 2 & 37 & 9  & 0.27 & 3 \\
geo+demand clustering$^{\dagger}$ & 2 & 29 & 4 & 0.76 & 0 \\
\textbf{SC-MAE-RG (proposed)} & 2 & \textbf{31} & \textbf{5} & 0.61 & \textbf{1} \\
\midrule
$K$-means (geo)               & 3 & 43 & 15 & 0.23 & 3 \\
Ward (geo)                    & 3 & 40 & 12 & 0.46 & 3 \\
geo+demand clustering$^{\dagger}$ & 3 & 36--39 & 9--11 & 0.54--0.67 & 2--3 \\
\textbf{SC-MAE-RG (proposed)} & 3 & \textbf{38} & \textbf{10} & 0.44 & 3 \\
\bottomrule
\end{tabular}

\vspace{2pt}
{\footnotesize $^{\dagger}$Both $K$-means and Ward isolate S001 into a
singleton group ($1{+}14$; $1{+}7{+}7$ / $1{+}5{+}9$) and are therefore not
operationally valid despite the low gap.}
\end{table}

\begin{table}[t]
\centering
\footnotesize
\setlength{\tabcolsep}{3pt}
\caption{Peak concurrent resource demand of the recommended $K{=}2$ partition,
against the available inventory. G1/G2 = task groups; the gap column is
$\max(0,\mathrm{G1}+\mathrm{G2}-\mathrm{inventory})$ per resource class
($D{=}31$, $G{=}5$).}
\label{tab:gap}
\begin{tabular}{lccccc}
\toprule
Resource class & G1 & G2 & Total & Inventory & Gap \\
\midrule
A airframes & 2 & 2 & 4 & 4 & 0 \\
A batteries & 2 & 2 & 4 & 6 & 0 \\
B airframes & 2 & 2 & 4 & 2 & 2 \\
B batteries & 3 & 3 & 6 & 4 & 2 \\
C airframes & 2 & 0 & 2 & 2 & 0 \\
C batteries & 4 & 0 & 4 & 4 & 0 \\
Relay airframes & 2 & 1 & 3 & 2 & 1 \\
Relay energy components & 3 & 1 & 4 & 6 & 0 \\
\bottomrule
\end{tabular}
\end{table}

\begin{table}[t]
\centering
\footnotesize
\setlength{\tabcolsep}{3pt}
\caption{Top five service areas of the entropy-weighted TOPSIS vulnerability
ranking (closeness coefficient; larger is more vulnerable).}
\label{tab:topsis}
\begin{tabular}{lccccc}
\toprule
Area & Mass (kg) & Volume (m$^3$) & Distance (m) & Closeness & Rank \\
\midrule
S001 & 154 & 0.394 & 3063 & 0.895 & 1 \\
S003 & 81 & 0.211 & 7261 & 0.441 & 2 \\
S002 & 81 & 0.211 & 7435 & 0.438 & 3 \\
S004 & 59 & 0.156 & 7746 & 0.274 & 4 \\
S006 & 59 & 0.156 & 3181 & 0.268 & 5 \\
\bottomrule
\end{tabular}
\end{table}

\subsection{Discussion}
The partitioning stage converts the coupled schedule into an actionable
decision: the recommended $K{=}2$ plan has a quantified gap of five resource
units (two B-type airframes, two B-type batteries and one relay airframe), so the
augmentation list is explicit rather than qualitative.

From a deployment perspective the outputs map directly onto field operations:
per-trip loading and route sheets executable without re-running any
optimization, four independently executable relay sorties spanning seven service
windows and verified point-by-point at one-second resolution, DEM-derived cruise
altitudes that keep the prescribed clearance, and a procurement list naming the
five missing units of the $K{=}2$ plan by type, where the zero-replication
control shows cross-group relay sharing, not vehicle count, to be binding.

Two threats to validity deserve emphasis. First, the energy model is validated
only against frozen computational artifacts: the independent recomputation
reproduces the relay-point count within 0.3\% and the relay energy within
1.6\%, but no hardware flight test has been performed, and wind, payload-drag
and battery-aging effects could shift the absolute figures. Second, relay
demand is identified on a five-second search grid over the link budget, which
underestimates the relay-required population relative to the one-second
verification grid; the 11683-sample coverage claim therefore holds for the
verified solution, while the candidate-hover search may miss short-duration
blockage events. Further limitations are that link feasibility is judged by a
binary DEM line-of-sight test with a fixed occlusion loss, that the environment
is deterministic and does not model weather or equipment faults. For Problem
2, optimality is certified for the scheduling sub-problem only on a fixed
trip set, where the exact CP-SAT benchmark quantifies the greedy scheduler gap
at about $15\%$ in makespan; the joint batching--routing--scheduling optimum is
not searched, so the $65.25$\,kWh main plan remains the best of the enumerated
candidate structures rather than a proven joint optimum. Because the parameters
are data-driven, the model chain can be transferred to other mountainous
disaster settings, but the reported fleet configuration and coverage figures
are instance-specific.

Transferring the chain to a new site requires re-running, not re-designing:
(i)~recompute the machine--area safe-payload table from a DEM of the region;
(ii)~survey the available fleet, batteries and relay components into the
parameter tables; (iii)~adopt the link parameters and re-identify the
relay-required population from the planned trajectories; and (iv)~run the four
stages in order, since each consumes only its predecessors' outputs. Only the
Problem-2 tour enumeration and the Problem-3 per-second verification grow
noticeably with instance size, and both are embarrassingly parallel.

% =====================================================================
\section{Conclusion}
% ---------------------------------------------------------------------
This paper built a four-stage ``transport scheduling--communication
support--resource allocation'' framework that runs on a single physical caliber,
from DEM-driven batching to partition-based resource accounting; its core stage
minimizes the relay resources that keep the delivery UAVs connected once their
delivery plan is fixed. The exact
batching model returns an 18-trip plan at 59.24\,kWh together with a verified
optimality benchmark; the heterogeneous routing model delivers all 80 boxes,
satisfies all 31 hard deadlines and recommends a 20-trip main plan at
65.25\,kWh; and the joint transport--communication stage covers all 11683
sampled relay-required points with two relay airframes and four sorties at a
joint energy of 70.68\,kWh. The partitioning stage then quantifies the resource gap and
produces an explicit augmentation list. The sensitivity studies delineate the
operating envelope of these plans: the batching solution tolerates return
margins up to $35\%$ before losing feasibility at $40\%$, the relay demand is
governed by the receiving threshold rather than by the terrain-clearance
margin, and the multi-hover mode is the only tested configuration that reaches
full coverage with the available two relay airframes.

Four directions merit further work. (i)~The one-second discrete verification
should be replaced by event-driven continuous-time link checking, so that
coverage guarantees no longer depend on a sampling grid and short-duration
blockage events cannot be missed. (ii)~The deterministic environment should be
extended with wind fields, weather disturbances and stochastic demand, turning
the single-scenario plans into robust or chance-constrained counterparts.
(iii)~The single-depot setting should be generalized to multiple dispatch
depots with cross-depot sharing of airframes and batteries, which would relax
the resource-replication penalty quantified here. (iv)~The scheduling layer should
be coupled with the batching-structure decision inside one exact or
adaptive large-neighborhood-search model---our CP-SAT benchmark certifies the
fixed-trip scheduling optimum but the joint optimum remains unsearched---and
the benchmark family should be extended to instances larger than 20 trips.

% =====================================================================
\section*{Acknowledgment}
This work originated from an entry to the China Postgraduate Mathematical
Contest in Modeling; the authors thank the contest organizers for providing the
mountainous test scenario and data.

\noindent\textbf{Funding:} This research received no specific grant from any
funding agency in the public, commercial, or not-for-profit sectors.

\noindent\textbf{Conflict of interest:} The authors declare no competing
interests.

\noindent\textbf{Data availability:} Code, test instance, benchmark suite and
verification pipeline:
\texttt{https://github.com/\allowbreak vscode-aiyaya/\allowbreak uav-flood-transport}
(archival DOI upon acceptance).

% =====================================================================
\begin{thebibliography}{99}
% NOTE: identifiers below were verified against Crossref / publisher pages /
%       CNKI-CQVIP records (2026-10-15). Obscure Chinese theses and trade-press
%       items carry no DOI in IEEE style; [in Chinese] marks the language.
\bibitem{cuu2026} V.~H.~Cuu, T.~L.~Huynh, and \v{Z}.~Koboevi\'c, ``UAV-assisted
MOSI/SOMI MIMO-FSO relay for resilient transport communication links,''
\emph{Automation}, vol.~7, no.~4, p.~107, 2026, doi: 10.3390/automation7040107.
\bibitem{wu2024} L.~Wu, D.~Liu, and Y.~Deng, ``Emergency logistics path
optimization based on UAV transportation,'' \emph{Logistics Sci-Tech}, vol.~47,
no.~18, pp.~86--90, 2024, doi: 10.13714/j.cnki.1002-3100.2024.18.016, [in Chinese].
\bibitem{cao2024} Y.~Chen, ``Radio-map-based cargo delivery trajectory design
for cellular-connected UAVs,'' M.S. thesis, Nanchang University, Nanchang,
China, 2024, [in Chinese].
\bibitem{wuXu2026} X.~Wu, ``Joint delivery routing optimization of UAV and ground
vehicles from the perspective of urban logistics,'' M.S. thesis, Civil Aviation
Flight University of China, Guanghan, China, 2026, [in Chinese].
\bibitem{wu2026} X.~Wu, M.~Wang, and H.~Hu, ``Joint delivery routing optimization
of UAV and ground vehicles under regional restrictions,'' \emph{Science
Technology and Industry}, vol.~26, no.~2, pp.~53--60, 2026, [in Chinese].
\bibitem{xia2019} H.~Xia, T.~Yan, and H.~Zhang, ``Research on mathematical
modeling algorithm for multi-objective UAV transportation network,'' \emph{Electro-Mechanical
Information}, vol.~20, pp.~87--89, 2019, [in Chinese].
\bibitem{sun2026} X.~Sun and D.~Zhang, ``Optimizing fruit transportation in
mountainous orchards: A mission planning method using cooperative drone-vehicle
systems,'' \emph{Smart Agricultural Technology}, vol.~14, p.~102383, 2026,
doi: 10.1016/j.atech.2026.102383.
\bibitem{dong2026} Z.~Dong, Y.~Qian, W.~Ye, W.~Li, and T.~Gu, ``Three-dimensional
path planning and scheduling of urban emergency UAVs under typhoon disasters,''
\emph{J. South China Univ. Technol. (Natural Science Edition)}, vol.~54, no.~6,
pp.~74--87, 2026, doi: 10.12141/j.issn.1000-565X.250421, [in Chinese].
\bibitem{wan2024} H.~Wan, ``Research on design methods of flood-search-and-rescue
UAVs,'' Master's thesis, Shanghai University of Engineering Science, Shanghai, China,
2024, [in Chinese].
\bibitem{sabuncu2024} \"O.~Sabuncu and B.~Bilgehan, ``Revolutionizing healthcare
5.0: Blockchain-driven optimization of drone-to-everything communication using 5G
network for enhanced medical services,'' \emph{Technology in Society}, vol.~77,
p.~102552, 2024, doi: 10.1016/j.techsoc.2024.102552.
\bibitem{huang2024} L.~Huang, D.~Qu, J.~Zhou, and J.~Zhang, ``The optimization of
UAV-assisted downlink transmission based on RSMA,'' \emph{Mathematics}, vol.~13,
no.~1, p.~13, 2025, doi: 10.3390/math13010013.
\bibitem{zhang2021} X.~Zhang, F.~Zhang, P.~Huang, J.~Gao, H.~Yu, C.~Pei, and
Y.~Zhang,
``Self-triggered based coordinate control with low communication for tethered
multi-UAV collaborative transportation,'' \emph{IEEE Robotics and Automation
Letters}, vol.~6, no.~2, pp.~1559--1566, 2021, doi: 10.1109/LRA.2021.3057294.
\bibitem{li2026} Y.-F.~Li, X.~Huang, C.-Y.~Guo, Y.-R.~Wang, S.-M.~Wu, and
J.-B.~Jian, ``UAV-UGV cooperative perception for low-altitude transportation:
A review,''
\emph{Acta Automatica Sinica}, vol.~52, no.~2, pp.~210--229, 2026,
doi: 10.16383/j.aas.c250386, [in Chinese].
\bibitem{ma2020} J.~Ma, ``Research on autonomous swarm algorithms of
heterogeneous multi-UAVs,'' M.S. thesis, Shenyang Aerospace University, Shenyang,
China, 2020, [in Chinese].
\bibitem{guo2026} Y.~Guo, P.~Gao, and Z.~Wu, ``Two-stage UAV facility location and
scheduling based on robust satisficing,'' \emph{Operations
Research Transactions}, vol.~30, no.~3, pp.~204--218, 2026, [in Chinese],
doi: 10.15960/j.cnki.issn.1007-6093.2026.03.012.
\bibitem{yanmaz2023} E.~Yanmaz, ``Joint or decoupled optimization: Multi-UAV path
planning for search and rescue,'' \emph{Ad Hoc Networks}, vol.~138, p.~103018,
2023, doi: 10.1016/j.adhoc.2022.103018.
\bibitem{soya2025} H.~Soya, K.~Inagaki, and H.~So, ``A flight route design method
considering multi-hop communication using delivery UAVs,'' \emph{Drones},
vol.~9, no.~11, p.~751, 2025, doi: 10.3390/drones9110751.
\end{thebibliography}

\end{document}
